\chapter{Register Read Stage}

\section{Overview}

Following register renaming, instructions enter the Register Read (RR) stage. At this point, all
architectural register references have already been converted into physical register identifiers.

The purpose of this stage is to obtain the actual operand values required for execution.

In Sargantana, the Register Read stage is centered around the physical register file.

RTL location:

\begin{center}
\rtl{rtl/datapath/rtl/rr\_stage/}
\end{center}

Primary implementation:

\begin{center}
\rtl{rtl/datapath/rtl/rr\_stage/rtl/regfile.sv}
\end{center}

This stage receives renamed instructions from the Rename stage and forwards operand values
toward the Execution stage.

\begin{figure}[h]
  \centering
  \includegraphics[width=0.5\textwidth]{ch08_register_read_overview.png}
  \caption{Register Read stage: RFiles, Bypass, and Collision Detection}
  \label{fig:register-read-overview}
\end{figure}

\section{Register File}

The register file stores all physical register values currently available inside the processor.
Unlike a simple in-order processor, the register file is accessed using physical register numbers
generated during register renaming.

The module is implemented in:

\begin{center}
\rtl{rtl/datapath/rtl/rr\_stage/rtl/regfile.sv}
\end{center}

The register file supports:
\begin{itemize}
  \item Reading source operands
  \item Writing completed results back into physical registers
  \item Supplying operands to multiple execution units
\end{itemize}

This module serves as the central storage element for operand values.

\section{Operand Read}

Once physical register identifiers are available, the register file reads the corresponding operand
values.

For an instruction:
\begin{verbatim}
ADD P17, P4, P8
\end{verbatim}

the Register Read stage performs:
\begin{verbatim}
Read P4
Read P8
  |
Operand A
Operand B
  |
Execution Stage
\end{verbatim}

This process occurs every cycle for all issued instructions.

The implementation is handled by:

\begin{center}
\rtl{rtl/datapath/rtl/rr\_stage/rtl/regfile.sv}
\end{center}

\section{Operand Forwarding}

Waiting for every instruction to write back before reading operands would significantly reduce
performance. To minimize stalls, recently computed values may be forwarded directly from
execution units.

\begin{figure}[h]
  \centering
  \includegraphics[width=\textwidth]{ch08_bypass_network.png}
  \caption{Conceptual bypass path: ALU Result $\rightarrow$ Bypass Network $\rightarrow$ Register Read Stage $\rightarrow$ Execution}
  \label{fig:bypass-network}
\end{figure}

The pipeline overview shows this Bypass path between the Register Read and Execution
stages.

Although the forwarding logic spans multiple execution modules, the Register Read stage is
designed to receive these forwarded values whenever available.

\section{Dependency Resolution}

Even after register renaming, true data dependencies (Read After Write) must still be respected.

The Register Read stage ensures that:
\begin{itemize}
  \item Ready operands are immediately read from the register file.
  \item Forwarded operands are used whenever newer values exist.
  \item Instructions with unresolved dependencies wait until their operands become available.
\end{itemize}

Once all operands are available, the instruction advances to the Execution stage, where
arithmetic, logical, memory, floating-point, SIMD, or branch operations are performed.

The interaction between Register Read and Execution is coordinated through:

\begin{center}
\rtl{rtl/datapath/rtl/datapath.sv}
\end{center}

and the execution subsystem rooted at:

\begin{center}
\rtl{rtl/datapath/rtl/exe\_stage/rtl/exe\_stage.sv}
\end{center}

This completes the front-end portion of the Sargantana datapath before execution begins.
